% ============================================================
% 60_boxes.tex — tcolorbox-Stile, Balken, Boxen, Formelnoten, Listen
% ============================================================

% ------------------------------------------------------------
% Balken-Bindung: ein Titelbalken steht nie allein am Spaltenfuss
% ------------------------------------------------------------
% Ein Balken ohne den Anfang seines Inhalts ist der auffälligste Layoutfehler
% einer ZSF: Man sieht die Überschrift, aber nicht, wozu sie gehört.
%
% Eine Platzreserve allein verhindert das nicht, sie verursacht es: Sie besteht
% aus \penalty0 … \penalty-100 und legt damit einen BILLIGEN Umbruchpunkt an
% genau die Stelle, an der ein Balken gerade fertig gesetzt wurde. Ihre Bremse
% — Stretch-Glue, die einen Umbruch bei viel Restplatz unattraktiv macht —
% trägt unter \raggedcolumns nicht, weil multicols den Spaltenboden ohnehin
% frei auffüllt. Übrig bliebe die blosse Einladung zum Umbruch.
%
% Deshalb zwei getrennte Aufgaben statt einer:
%   \ZSFReserveSpace  — „lieber vorher umbrechen, wenn es eng wird" (Bremse)
%   \ZSFTitleBarLock  — „hier NICHT umbrechen, es kam noch kein Inhalt" (Sperre)
% Die Sperre ist die verlässliche; sie hängt an keiner Höhenschätzung.
\newif\ifzsfTitleBarPending
\newcommand{\ZSFTitleBarLockOn}{\global\zsfTitleBarPendingtrue}
\newcommand{\ZSFTitleBarLockOff}{\global\zsfTitleBarPendingfalse}

% Jeder Umbruchpunkt, den das Balken-System selbst setzt, läuft hierüber:
% solange noch kein Inhalt folgte, wird daraus ein Verbot.
\newcommand{\ZSFBarPenalty}[1]{%
  \ifzsfTitleBarPending\penalty10000\else\penalty#1\fi
}

\newcommand{\ZSFReserveSpace}[1]{%
  \par
  \ifzsfTitleBarPending
    \penalty10000
  \else
    \vspace{#1}\penalty0\vspace{-#1}%
  \fi
  \relax
}
\newcommand{\BreakIfNotEnoughSpace}{\ZSFReserveSpace{\ZSFbreakThreshold}}
\newcommand{\BreakIfNotEnoughSpaceChapter}{\ZSFReserveSpace{\ZSFbreakThresholdChapter}}
\newcommand{\BreakIfNotEnoughSpaceSubsubsection}{\ZSFReserveSpace{\ZSFbreakThresholdSubsubsection}}

% Einheitlicher Top-Spacing-Kontrakt für Titelbalken.
% Reihenfolge ist kritisch: \addvspace MUSS vor \BreakIfNotEnoughSpace stehen,
% damit es per Max-Semantik mit dem after-skip der Vorbox kollabiert.
% tcolorbox-Definitionen, die diesen Helper rufen, MUESSEN before skip=0pt
% setzen, damit das Spacing nicht doppelt eingerechnet wird.
\newcommand{\ZSFTitleBarTopSpacing}[1]{%
  \par
  \addvspace{#1}%
  \BreakIfNotEnoughSpace
}

%
%
%
%
%
%
%
%
%
%
%
%
%
%
%
%
%
%
%
%
%
%
%
%
%
%
\newcommand{\ZSFTitleBarAfter}{%
  \par\penalty10000\ZSFInterlude{\ZSFbarAfterGap}\ZSFTitleBarLockOn
}

\tcbset{
  zsfbarcolor/.style={colback=#1, colframe=#1},
  zsfbar/.style={
    enhanced jigsaw,
    breakable=false,
    boxrule=0pt,
    boxsep=0pt,
    arc=\ZSFboxArc,
    outer arc=\ZSFboxArc,
    left=\ZSFboxPadX,
    right=\ZSFboxPadX,
    before skip=0pt,
  },
}

\newtcolorbox{chapterbar}[1][]{
  zsfbar,
  zsfbarcolor=\ChapterBarColor,
  top=\ZSFbarPadY,
  bottom=\ZSFbarPadY,
  before skip=\ZSFspaceL,
  after={\ZSFTitleBarAfter},
  fontupper=\ZSFfontChapter,
  #1
}
% \ZSFInkOwned steht in ALLEN vier Balken direkt neben der Kontrastfarbe:
% Ein Balken ist eine gesaettigte Flaeche mit eigener Textfarbe und besitzt
% damit die Farbe (40_colors_structure -> Ink-Vertrag). Bewusst gleich
% behandelt, auch wo die Flaeche hell ist (Unterabschnitt) — das Kriterium
% ist "setzt eigene Kontrastfarbe" und nicht eine Einzelfallabwaegung, sonst
% muesste man es bei jeder Farbaenderung neu entscheiden.
\newcommand{\ChapterBar}[1]{
  \BreakIfNotEnoughSpaceChapter
  \begin{chapterbar}
    \textcolor{\ChapterBarTextColor}{\ZSFInkOwned #1}
  \end{chapterbar}
}

% ------------------------------------------------------------
% Dokument-Titel-Masthead — einmalig vor dem Front-Chapter
% ------------------------------------------------------------
\newtcolorbox{zsftitleheader}[1][]{
  zsfbar,
  zsfbarcolor=ZSFDocumentTitleBarColor,
  left=\ZSFboxPadXBar,
  right=\ZSFboxPadXBar,
  top=\ZSFbarPadY,
  bottom=\ZSFbarPadY,
  after={\ZSFTitleBarAfter},
  halign=center,
  #1
}
\newcommand{\ZSFTitleHeader}[2]{%
  \ZSFBarPenalty{-500}%
  \begin{zsftitleheader}
    {\color{ZSFDocumentTitleTextColor}\ZSFInkOwned\ZSFfontTitleHeader #1\par}%
    \vspace{\ZSFspaceXS}%
    {\color{ZSFDocumentTitleBylineColor}\ZSFInkOwned\ZSFfontNote von~#2\par}%
  \end{zsftitleheader}
}

% ------------------------------------------------------------
% Subsection-Bar (+ unnummerierte Star-Variante)
% ------------------------------------------------------------
\newtcolorbox{subsectionbar}[1][]{
  zsfbar,
  zsfbarcolor=\SubsectionBarColor,
  top=\ZSFbarPadY,
  bottom=\ZSFbarPadY,
  after={\ZSFTitleBarAfter},
  fontupper=\ZSFfontSection,
  #1
}
\newtcolorbox{subsubsectionbar}[1][]{
  zsfbar,
  zsfbarcolor=\SubsubsectionBarColor,
  top=\ZSFbarPadYTight,
  bottom=\ZSFbarPadYTight,
  after={\ZSFTitleBarAfter},
  fontupper=\ZSFfontSubsubsection,
  #1
}
\makeatletter
% Pro Ebene EIN Rumpf; die Varianten liefern nur das Nummernpräfix (leer bei
% der Sternvariante).
%   #1 = Label (leer erlaubt)   #2 = Nummernpräfix   #3 = Titel
\newcommand{\ZSF@subsectionBarBody}[3]{%
  \ZSFTitleBarTopSpacing{\ZSFsubsectionBarBeforeSkip}%
  \ZSFBarPenalty{-500}%
  \ZSFCollapseAfterChapterBar
  \ZSF@storeLabelName{#3}%
  \if\relax\detokenize{#1}\relax\else\label{#1}\fi
  \begin{subsectionbar}
    \textcolor{\SubsectionBarTextColor}{\ZSFInkOwned #2#3}
  \end{subsectionbar}
}
\newcommand{\ZSF@subsubsectionBarBody}[3]{%
  \BreakIfNotEnoughSpaceSubsubsection
  \par\addvspace{\ZSFsubsubsectionBarBeforeSkip}%
  \ZSFBarPenalty{-300}%
  \ZSF@storeLabelName{#3}%
  \if\relax\detokenize{#1}\relax\else\label{#1}\fi
  \begin{subsubsectionbar}
    \textcolor{\SubsubsectionTitleColor}{\ZSFInkOwned #2#3}
  \end{subsubsectionbar}
}

\newcommand{\SubsectionBar}{\@ifstar{\SubsectionBarStar}{\SubsectionBarNumbered}}
\newcommand{\SubsectionBarNumbered}[2][]{%
  \refstepcounter{subsection}\setcounter{subsubsection}{0}%
  \ZSF@subsectionBarBody{#1}{\thesubsection. }{#2}}
\newcommand{\SubsectionBarStar}[2][]{\ZSF@subsectionBarBody{#1}{}{#2}}

\newcommand{\SubsubsectionBar}{\@ifstar{\SubsubsectionBarStar}{\SubsubsectionBarNumbered}}
\newcommand{\SubsubsectionBarNumbered}[2][]{%
  \refstepcounter{subsubsection}%
  \ZSF@subsubsectionBarBody{#1}{\thesubsubsection. }{#2}}
\newcommand{\SubsubsectionBarStar}[2][]{\ZSF@subsubsectionBarBody{#1}{}{#2}}

\makeatother

%
%
%
%
\newcommand{\ZSFNewColumn}{\columnbreak}

%
%
%
%
%
%
%
%
%
%
%
%
%
%
%
%
%
%
%
%
%
%
%
%
%
%
%
%
%
%
%
\tcbset{zsftitlebar/.style={
  fonttitle=\ZSFfontBoxTitle\ZSFInkOwned\strut,
  halign title=center,
  title style={minimum height=\ZSFtitleBarHeight},
  toptitle=\ZSFspaceXS,
  bottomtitle=\ZSFspaceXS,
  before title={},
  after title={},
}}

% Rechtsbuendiges Meta-/Kategorie-Tag im Titelbalken einer Titelbox.
% Verwendung im Titel-Argument der Box, z.B.:
%   \begin{defbox}[Quicksort\ZSFTitleTag{O(n log n)}] ... \end{defbox}
%   \begin{codebox}[Algorithmus\ZSFTitleTag{Python}]  ... \end{codebox}
% Schiebt das Tag per \hfill an den rechten Titelrand. Schrift: \ZSFfontTitleTag.
\newcommand{\ZSFTitleTag}[1]{\hfill{\ZSFfontTitleTag #1}}

% Zum "before"-Schlüssel in allen Box-Stilen unten: das führende \par ist Pflicht.
% "before" ERSETZT die Default-Aktion von tcolorbox (die selbst ein \par enthält) —
% es hängt nicht an. Ohne \par wird eine Box, vor der roher Fliesstext ohne
% Leerzeile steht, inline an den offenen Absatz gehängt: sie bekommt keine
% Spaltenbreite und ragt sichtbar in die Nachbarspalte. Im vertikalen Modus ist
% \par ein No-op, der Zusatz ist also folgenlos, wo ohnehin alles korrekt war.
% Gemeinsamer Grundvertrag aller Inhaltsboxen. Was hier steht, gilt fuer jede
% Box — die einzelnen Stile setzen nur noch ihre Unterschiede. Besonders das
% before-Hook ist Pflicht (Begruendung im Kommentar oben) und darf deshalb
% nicht pro Box wiederholt werden, wo man es vergessen kann.
% Schicht 1 — Fluss. Auch rahmenlose Boxen (weight=quiet, splitbox) laufen
% hierüber, damit Abstände und Balken-Bindung überall identisch greifen.
% DER ABSTAND VOR EINEM BLOCK — eine Stelle, und sie muss es sein.
%
% Hier standen `before skip` und `before` nebeneinander. In tcolorbox ist
% `before skip` aber KEIN eigener Schluessel, sondern ein Stil, der `before`
% schreibt (tcolorbox.sty: `before skip/.style={before={...}}`). Beide
% schreiben also dasselbe Makro, und der spaetere gewinnt — `before` stand
% hinten, damit war \ZSFBoxBeforeSkip wirkungslos. Nachgemessen: auf 40pt
% gesetzt aendert sich am Satz nichts.
%
% Die Folgen reichten weiter als der eine Wert: Der Abstand vor JEDER Box kam
% damit allein aus \parskip und gehoerte keinem Register. Die Kollaps-Mechanik
% nach einem Balken (damals \ifzsfAfterSubsectionBar) rechnete einen Wert aus, den
% niemand las, und \ZSFboxgroup schloss seine Zwischenraeume nicht wirklich.
%
% Deshalb EIN `before`, das beides tut. Die Skip-Logik ist die von tcolorbox
% selbst — dieselben Faelle (nichts auf der Liste, vertikaler Modus), damit
% hier keine zweite Auslegung derselben Frage entsteht. Gelesen wird
% \ZSFBoxBeforeSkip VOR dem Verbrauch des Merkers: Der Merker waehlt den Wert.
%
% Den Minipage-Sonderfall von tcolorbox braucht es nicht: \addvspace prueft
% LaTeXs eigenes @minipage selbst und tut am Anfang einer Minipage nichts —
% genau das, was jener Zweig von Hand herstellt. Ihn nachzubauen hiesse,
% \iftcb@minipage zu lesen, und das existiert ausserhalb einer tcolorbox gar
% nicht (gemessen: undefined im runintext).
\makeatletter
\newcommand{\ZSF@blockBefore}{%
  \ifnum\lastnodetype=-1\relax\else
    \par
    \ifvmode\addvspace{\glueexpr\ZSFBoxBeforeSkip-\parskip\relax}\fi
  \fi
  % BEIDE Merker, nicht nur der eine. Ein Block zwischen Kapitelbalken und
  % Abschnittsbalken heisst, dass die beiden Balken NICHT aneinanderstossen —
  % blieb der Kapitel-Merker stehen, zog der Abschnittsbalken seinen Vorabstand
  % trotzdem zusammen. Derselbe Fehler eine Ebene hoeher.
  \ZSFConsumeAfterChapterBar\ZSFTitleBarLockOff
}
\tcbset{zsfbox/.style={
  enhanced jigsaw,
  after skip=\ZSFBoxAfterSkip,
  before={\ZSF@blockBefore},
}}
\makeatother

%
%
%
%
%
%
%
%
%
%
%

% Polsterung ist zweidimensional: waagrecht entscheidet über Kante und
% Balkenfreiraum, senkrecht über die Höhe. Getrennt, weil die Fälle sich
% tatsächlich kreuzen (Wertegitter: waagrecht bündig, senkrecht knapp).
%
% `padx` setzt den Abstand von Inhalt UND Titel zur Kante — es ist eine Kante
% und deshalb ein Wert (Begründung bei zsftitlebar oben). Für zwei der drei
% Werte steht deshalb NUR `left`/`right` da: Das sind Sammelschlüssel von
% tcolorbox (`left={lefttitle=#1, leftupper=#1, leftlower=#1}`), der Titel
% folgt also von selbst. Ihn danebenzuschreiben wäre derselbe Wert zweimal —
% und beim nächsten Eingriff einmal anders.
%
% `padx=none` ist der Fall, in dem Titel und Inhalt auseinandergehen MÜSSEN,
% und deshalb der einzige, der beide nennt: Die Box polstert nicht, weil ihr
% Inhalt es selbst tut — tablebox und valuegrid lassen ihre Farbfläche bis an
% den Rahmen laufen und halten den Rand in der äussersten Zelle
% (\ZSFtableEdgePad, siehe 20_tables). Der Titel hat keine Zelle und muss den
% Wert direkt lesen; ohne diese Zeile stünde er als einziges Element der Box
% an der blanken Kante.
\tcbset{
  padx/.is choice,
  padx/normal/.style={left=\ZSFboxPadX,    right=\ZSFboxPadX},
  padx/bar/.style   ={left=\ZSFboxPadXBar, right=\ZSFboxPadX},
  padx/none/.style  ={left=0pt,            right=0pt,
                      lefttitle=\ZSFtableEdgePad, righttitle=\ZSFtableEdgePad},
  pady/.is choice,
  pady/normal/.style={top=\ZSFboxPadY,    bottom=\ZSFboxPadY},
  pady/tight/.style ={top=\ZSFtextMinPad, bottom=\ZSFtextMinPad},
  pady/none/.style  ={top=0pt,            bottom=0pt},
}

%
%
%
%
%
%
%
%
%
%
%
%
%
%
%
%
\newenvironment{ZSFinset}{%
  \par\advance\leftskip\ZSFboxPadX\advance\rightskip\ZSFboxPadX
}{%
  \par
}

% Ausrichtung des Boxinhalts. Als Hook statt als direktem 'before upper',
% damit die Eintritts-Sequenz einer Box an EINER Stelle steht und die
% Ausrichtung sie nicht ersetzen muss.
\newcommand{\ZSFBoxAlign}{}
% Kontext-Flag der Box (heute nur die formulabox). Ebenfalls Teil derselben
% Sequenz — sonst bräuchte jede Box mit Kontext ein eigenes 'before upper'.
\newcommand{\ZSFBoxContext}{}
% Absatzabstand im Boxinhalt. Textboxen setzen ihn eigenständig (\ZSFspaceS),
% Container um Tabellen und Listen erben den Dokumentwert — dort gibt es keine
% Absätze, und ein eigener Wert würde nur die Liste verschieben. Als Regler
% statt als Auslassung, damit „erbt bewusst" und „vergessen" unterscheidbar sind.
\newcommand{\ZSFBoxParskip}{}
%
%
%
%
%
%
%
%
%
\newcommand{\ZSFBoxFont}{}
% Die Ausrichtung ERSETZT die Justierung, sie ergänzt sie nicht: `align=left`
% ist der Flattersatz für schmale Spalten (55_readability), `align=center` ist
% zentrierter Satz. Beides zusammen wäre kein zentrierter Flattersatz, sondern
% zentrierter Satz mit den Umbruch-Penalties des Flattersatzes — sichtbar an
% anders gebrochenen Formelzeilen.
\tcbset{
  align/.is choice,
  align/left/.style  ={/utils/exec={\renewcommand{\ZSFBoxAlign}{\ZSFReadableAuto}}},
  align/center/.style={/utils/exec={\renewcommand{\ZSFBoxAlign}{\centering}}},
  bodyparskip/.is choice,
  bodyparskip/box/.style={
    /utils/exec={\renewcommand{\ZSFBoxParskip}{\setlength{\parskip}{\ZSFspaceS}}}},
  bodyparskip/inherit/.style={
    /utils/exec={\renewcommand{\ZSFBoxParskip}{}}},
  font/.is choice,
  font/normal/.style={/utils/exec={\renewcommand{\ZSFBoxFont}{\ZSFfontContent}}},
  font/dense/.style ={/utils/exec={\renewcommand{\ZSFBoxFont}{\ZSFfontContentDense}}},
}
% Titel-Merker der leisen Fassung. Gefüllt wird er von der Fabrik, gesetzt vom
% Gewicht-Regler weiter unten; hier stehen nur die Deklarationen, damit sie vor
% beidem existieren.
%
% \ifZSF@boxQuiet trägt die Gewicht-Wahl. Sie steht schon fest, BEVOR die
% Optionsliste läuft — die Fabrik liest sie in einem Vor-Lauf aus dem Aufruf
% (siehe dort). Deshalb darf der Ton weiter unten sie bereits abfragen.
\makeatletter
\newcommand{\ZSF@boxTitleText}{}
\newif\ifZSF@boxHasTitle
\newif\ifZSF@boxQuiet
\makeatother
\newcommand{\ZSFBoxQuietTitle}{}
%
%
%
%
%
%
%
%
%
%
%
%
%
%
%
\newcommand{\ZSFBoxSetup}{\ZSFBoxFont\ZSFBoxParskip\ZSFBoxAlign\ZSFBoxContext}
\newcommand{\ZSFBoxEnter}{\ZSFBoxSetup\ZSFBoxQuietTitle}

%
%
%
%
%
%
%
%
%
%
%
%
%
%
%
%
%
%
\tcbset{zsfboxcontract/.style={
  zsfbox,
  tone=chapter,
  align=left,
  bodyparskip=box,
  font=normal,
  % Hooks zurücksetzen, damit eine verschachtelte Box die Einstellungen der
  % äusseren nicht erbt — sie sind gruppenlokal, aber nicht selbstlöschend.
  /utils/exec={\renewcommand{\ZSFBoxContext}{}\renewcommand{\ZSFBoxQuietTitle}{}},
  before upper={\ZSFBoxEnter},
  % Der Abstand zwischen oberer und unterer Haelfte, wenn sie gestapelt
  % stehen. Er gehoert in den Vertrag und nicht in ein Preset: JEDE Box kann
  % ein \tcblower tragen. Gesetzt hatte ihn nur die formulabox — alle
  % uebrigen bekamen tcolorbox' 2mm, ein Mass, das niemand gewaehlt hat und
  % das der Dichte nicht folgt. Derselbe Wert wie \ZSFsep: Es ist derselbe
  % Blockwechsel, nur von tcolorbox gezogen statt vom Autor.
  middle=\ZSFspaceSep,
  frame=soft,
  ZSF@surface=auto,
  ZSF@titlefill=tone,
  splitalign=top,
}}

%
%
%
%
%
%
%
%
%
%
%
%
%
%
%
%
%
%
%
%
%
%
\makeatletter
\newcommand{\ZSF@frameChoice}{soft}
\tcbset{
  frame/.is choice,
  frame/soft/.code  ={\renewcommand{\ZSF@frameChoice}{soft}},
  frame/strong/.code={\renewcommand{\ZSF@frameChoice}{strong}},
  frame/hard/.code  ={\renewcommand{\ZSF@frameChoice}{hard}},
  frame/none/.code  ={\renewcommand{\ZSF@frameChoice}{none}},
}
% Läuft als letzte Option jeder Box und übersetzt die gemerkte Stärke gegen
% den dann feststehenden Ton. Jede Box, die einen Rahmen haben kann, muss ihn
% durchlaufen — sonst bleibt sie beim tcolorbox-Default.
\tcbset{ZSF@resolveframe/.code={%
  \IfStrEqCase{\ZSF@frameChoice}{%
    {soft}  {\tcbset{colframe=\ZSFtoneFrameSoft}}%
    {strong}{\tcbset{colframe=\ZSFtoneFrameStrong}}%
    {hard}  {\tcbset{colframe=\ZSFtoneFrameHard}}%
    {none}  {\tcbset{frame hidden, boxrule=0pt, arc=0pt, outer arc=0pt}}%
  }%
}}
\makeatother

% Fläche — der Baustein wählt die ROLLE, die Farbe dazu kommt aus dem Ton.
%
% Wortgleich zum Rahmen darüber und aus demselben Grund: tcolorbox löst
% `colback` und `colbacktitle` beim Setzen der Option sofort per \colorlet auf.
% Beide gehören damit dem letzten Schreiber. Ein Preset, das seine Fläche
% direkt setzt, verliert sie deshalb an jede spätere Tonwahl des Aufrufers —
% und ein ausgeschriebenes `tone=chapter` sähe anders aus als ein
% weggelassenes, obwohl es dieselbe Wahl ist. Verzögert aufgelöst stehen
% Rolle und Ton beide fest, wenn der Resolver läuft; die Reihenfolge im
% Aufruf ist dann bedeutungslos.
%
%   auto      laut oder ruhig — entscheidet das Gewicht (Default)
%   plain     ungetönt; Basis für Zebra, Gitterlinien, Syntaxfarben
%   quiet     immer die ruhige Fläche, unabhängig vom Gewicht
%   emphasis  betont; für Boxen, deren Inhalt selbst die Aussage ist
%
% `ZSF@titlefill` beantwortet dieselbe Frage für die Kopfzeile: `tone` gibt
% ihr den Titelbalken des Tons, `surface` setzt sie auf die Fläche der Box
% (Formelbox — dort ist der Titel eine Zeile, kein Balken).
\makeatletter
\newcommand{\ZSF@surfaceRole}{auto}
\newcommand{\ZSF@titleFillRole}{tone}
\tcbset{
  ZSF@surface/.is choice,
  ZSF@surface/auto/.code    ={\renewcommand{\ZSF@surfaceRole}{auto}},
  ZSF@surface/plain/.code   ={\renewcommand{\ZSF@surfaceRole}{plain}},
  ZSF@surface/quiet/.code   ={\renewcommand{\ZSF@surfaceRole}{quiet}},
  ZSF@surface/emphasis/.code={\renewcommand{\ZSF@surfaceRole}{emphasis}},
  ZSF@titlefill/.is choice,
  ZSF@titlefill/tone/.code   ={\renewcommand{\ZSF@titleFillRole}{tone}},
  ZSF@titlefill/surface/.code={\renewcommand{\ZSF@titleFillRole}{surface}},
}
% Läuft als letzte Option jeder Box, direkt neben dem Rahmen-Resolver.
% Erst die Rolle in eine Farbe übersetzen, dann setzen — so kann die Kopfzeile
% dieselbe Farbe bekommen, ohne die Fallunterscheidung zu wiederholen.
\newcommand{\ZSF@resolvedBack}{\ZSFtoneLoudBack}
\tcbset{ZSF@resolvesurface/.code={%
  \IfStrEqCase{\ZSF@surfaceRole}{%
    {auto}    {\ifZSF@boxQuiet\renewcommand{\ZSF@resolvedBack}{\ZSFtoneQuietBack}%
               \else\renewcommand{\ZSF@resolvedBack}{\ZSFtoneLoudBack}\fi}%
    {plain}   {\renewcommand{\ZSF@resolvedBack}{\BoxPlainBackColor}}%
    {quiet}   {\renewcommand{\ZSF@resolvedBack}{\ZSFtoneQuietBack}}%
    {emphasis}{\renewcommand{\ZSF@resolvedBack}{\ZSFtoneEmphasisBack}}%
  }%
  \tcbset{colback=\ZSF@resolvedBack}%
  \IfStrEqCase{\ZSF@titleFillRole}{%
    {tone}   {}%
    {surface}{\tcbset{colbacktitle=\ZSF@resolvedBack}}%
  }%
}}
\makeatother

% Atomar: bricht nie über eine Spaltengrenze, auch nicht im Pack-Modus.
% Muss in der Optionsliste NACH dem Basisstil stehen — dorthin hängt
% \ZSFBoxesBreakableOn sein 'breakable'. Als benannter Regler ist diese
% Reihenfolge in den Presets festgeschrieben statt in einem Kommentar.
\tcbset{atomic/.style={unbreakable}}

% Schicht 2 — Form. Alles, was die Geometrie einer geschlossenen Inhaltsbox
% bestimmt. Wer Rahmenstärke, Ecken oder Innenabstand des Systems ändern will,
% ändert hier — und nur hier. padx/pady gehören dazu, damit \linewidth in jeder
% Box identisch ist (sonst variierende Tabellen-Breiten je Box-Typ).
%
% Was NICHT hier steht, sondern im Vertrag darüber: alles, was auch für eine
% Box mit eigener Form gilt (Ton, Justierung, Eintritt, Reset der verzögerten
% Wahlen). Sonst müsste jede rahmenlose Box es abschreiben.
\tcbset{zsfboxshape/.style={
  zsfboxcontract,
  breakable=false,
  boxsep=0pt,
  boxrule=\ZSFboxRuleWidth,
  titlerule=0pt,
  arc=\ZSFboxInnerArc,
  outer arc=\ZSFboxArc,
  padx=normal,
  pady=normal,
}}

%
%
%
%
%
%
%
%
%
%
%
%
%
%
%
%
%
%
%
%
%
%
%
%
%
%
%
%
%
%
%
%
%
%
%
%
%
%
%
%
%
%

% Schicht 3 — Titelbox. Der EINZIGE Vertrag für Boxen mit Titelbalken; alle
% Presets unten stehen darauf.
%
% Dass hier nur noch zwei Zeilen stehen, ist der Punkt: Form kommt aus
% zsfboxshape, alles Übrige aus dem Vertrag darüber. Eine Titelbox ist damit
% genau das, was ihr Name sagt — eine Box mit Titelbalken.
\tcbset{zsftitlebox/.style={
  zsfboxshape,
  zsftitlebar,
}}

%
%
%
%
%
%
%
%
%
\makeatletter
% Die Fallunterscheidung sitzt in einem tcolorbox-Schlüssel, nicht um das
% \begin herum: So gibt es genau EIN \begin{tcolorbox} und genau ein \end —
% sonst zählt chktex sie gegeneinander auf und meldet sie bis zum Dateiende
% als unbalanciert (Warnung 15, dort nicht mehr lokal unterdrückbar).
\tcbset{zsftitle/.code={\IfBlankTF{#1}{\tcbset{notitle}}{\tcbset{title={#1}}}}}
%
%
%
%
%
%
%
%
%
%
%
%
%
%
%
%
%
%
%
%
%
%
%
%
%
%
%
%
%
%
%
%
%
\newcommand{\ZSF@ownedKeyError}[2]{%
  \PackageError{ZSF}{%
    Schluessel '#1' gehoert dem Box-System, nicht dem Aufruf%
  }{%
    Flaeche und Rahmen werden am Ende jeder Optionsliste aus Rolle und Ton
    zusammengesetzt (ZSF@resolveframe / ZSF@resolvesurface). Ein '#1' im
    Aufruf wird dabei ueberschrieben und bliebe wirkungslos -- frueher
    stillschweigend, jetzt hier. Stattdessen den Regler '#2' verwenden;
    fehlt dafuer ein Wert, ist das eine Luecke und gehoert gemeldet
    (fehlender Optionswert).%
  }%
}
% Box-spezifische Regler gehören genau einer Box. Auf jeder anderen nimmt
% tcolorbox sie klaglos entgegen, und niemand liest sie je wieder — dieselbe
% stille Klasse wie die Farbschlüssel oben, nur eine Ecke weiter:
% `\begin{defbox}[T][ordered]` setzte die Listen-Nummerierung, ohne dass eine
% Liste existiert; `[grid=none]` leerte die Gitterwahl einer Box ohne Gitter.
%
% Welche Box welchen Regler besitzt, sagt sie selbst, bevor sie die Fabrik
% ruft; leer heisst »keinen«. Die Kommas sind Begrenzer, damit ein Reglername,
% der Teil eines anderen ist, nicht versehentlich durchgeht.
\newcommand{\ZSF@boxOwnKnobs}{}
\newcommand{\ZSF@requireOwnKnob}[2]{%
  \IfSubStr{\ZSF@boxOwnKnobs}{,#1,}{}{%
    \PackageError{ZSF}{%
      Regler '#1' gehoert zu #2, nicht zu dieser Box%
    }{%
      Auf dieser Box hat '#1' nichts, worauf er wirken koennte -- er wuerde
      stillschweigend verschwinden. Entweder die Box verwenden, der er
      gehoert (#2), oder die Option weglassen.%
    }%
  }%
}
\pgfkeys{
  /zsf/boxpre/.is family, /zsf/boxpre,
  weight/.is choice,
  weight/loud/.code ={\ZSF@boxQuietfalse},
  weight/quiet/.code={\ZSF@boxQuiettrue},
  ordered/.code={\ZSF@requireOwnKnob{ordered}{ZSFlist}},
  grid/.code   ={\ZSF@requireOwnKnob{grid}{valuegrid}},
  part/.code   ={\ZSF@requireOwnKnob{part}{codebox}},
%
%
%
%
%
%
  colback/.code     ={\ZSF@ownedKeyError{colback}{tone}},
  colframe/.code    ={\ZSF@ownedKeyError{colframe}{frame}},
  colbacktitle/.code={\ZSF@ownedKeyError{colbacktitle}{tone}},
  coltitle/.code    ={\ZSF@ownedKeyError{coltitle}{tone}},
  .unknown/.code={},% chktex 26
}
\NewDocumentEnvironment{ZSF@box}{m m m}{%
  \renewcommand{\ZSF@boxTitleText}{#2}%
  \IfBlankTF{#2}{\ZSF@boxHasTitlefalse}{\ZSF@boxHasTitletrue}%
  % Zentraler Reset statt einer Zeile pro Basisstil: Sonst erbt eine Box, deren
  % Stil den Merker nicht zurücksetzt, die Titelzeile der Box, in der sie steht.
  \renewcommand{\ZSFBoxQuietTitle}{}%
  \ZSF@boxQuietfalse
  \pgfkeys{/zsf/boxpre, #3}%
  % Der Besitz gilt für GENAU diese Box. Zurückgesetzt direkt nach dem
  % Vor-Lauf, damit eine verschachtelte Box ihn nicht erbt und dort ein
  % fremder Regler wieder stumm durchginge.
  \renewcommand{\ZSF@boxOwnKnobs}{}%
  \begin{tcolorbox}[#1, zsftitle={#2}, ZSF@weightBase, #3,
                    ZSF@resolveframe, ZSF@resolvesurface, ZSF@resolvesplit]%
}{%
  \end{tcolorbox}%
}
\makeatother

% ------------------------------------------------------------
% Presets — jede Box ist eine Liste aus den Reglern oben
% ------------------------------------------------------------
% Prüffrage an jedes Preset: Steht hier ein Schlüssel, für den es einen Regler
% gäbe? Dann gehört der Regler hin. Rohe Schlüssel bleiben nur, wo die
% Eigenschaft wirklich einmalig ist (der Zebra-Beschnitt der Tabelle, die harte
% Kontur der Formelbox) — und dann mit Begründung.

% Tabelle: Hintergrund bis an die Kante, damit das Zebra nicht kurz vor dem
% Rahmen endet. 'clip upper' schneidet die Streifen an der Rundung ab.
% Die Fläche bleibt ungetönt, weil das Zebra sie als Basis braucht.
\tcbset{preset/table/.style={
  zsftitlebox,
  ZSF@surface=plain,
  padx=none, pady=none,
  bodyparskip=inherit,
  clip upper,
}}
%
%
\tcbset{preset/fig/.style={zsftitlebox, atomic}}

% ------------------------------------------------------------
% Gewicht — laut oder leise, dieselbe Box
% ------------------------------------------------------------
% Laut oder leise ist eine Eigenschaft derselben Box, keine zweite Box:
%
%   weight=loud   Titelbalken, Rahmen, getönte Fläche  (Default)
%   weight=quiet  linker Akzentbalken, rahmenlos, ruhige Fläche
%
% Der Titel wandert dabei aus der Kopfzeile in die erste Inhaltszeile. Den
% Titeltext liefert die Fabrik (siehe dort), deshalb gilt der Regler auf JEDER
% Box und nicht nur auf der defbox.
\makeatletter
%
%
%
%
%
%
%
%
%
\newcommand{\ZSF@emitQuietTitle}{%
  \ifZSF@boxHasTitle
    {\ZSFfontBoxTitle\color{\ZSFtoneAccent}\ZSFInkOwned\ZSF@boxTitleText}%
    \par\vspace{\ZSFspaceXS}%
  \fi
}
% Die Basis der leisen Fassung. Sie wird von der Fabrik VOR den
% Aufrufer-Optionen eingesetzt (Begründung dort), nicht mittendrin — deshalb
% darf sie Vorbelegungen wie `frame=none` oder `padx=bar` mitbringen, ohne
% eine Wahl des Aufrufers zu verwerfen.
%
% Die Fläche setzt sie bewusst NICHT selbst: Sie kommt vom Ton, der beide
% Fassungen kennt und anhand des Gewichts wählt (siehe `ZSF@toneSurface`).
% Täte sie es hier, wäre sie wieder ein Schreiber auf `colback` neben dem Ton
% — und welcher gewinnt, hinge erneut an der Reihenfolge.
\tcbset{ZSF@weightBase/.code={%
  \ifZSF@boxQuiet
    \tcbset{preset/quiet, notitle,
            /utils/exec={\renewcommand{\ZSFBoxQuietTitle}{\ZSF@emitQuietTitle}}}%
  \fi
}}
% Der Schlüssel selbst bleibt als No-op bestehen: Er steht weiterhin in der
% Aufrufer-Liste, wo tcolorbox ihn sonst als unbekannt meldete. Gelesen wurde
% er zu diesem Zeitpunkt längst — im Vor-Lauf der Fabrik.
\tcbset{
  weight/.is choice,
  weight/loud/.code ={},
  weight/quiet/.code={},
}
\makeatother

\NewDocumentEnvironment{defbox}{O{} O{}}
  {\begin{ZSF@box}{zsftitlebox}{#1}{#2}}{\end{ZSF@box}} % chktex 9

\NewDocumentEnvironment{tablebox}{O{} O{}}
  {\begin{ZSF@box}{preset/table}{#1}{#2}}{\end{ZSF@box}} % chktex 9
% figbox und splitbox tragen einen Block — dort gilt fuer ein \ZSFimage das
% Budget einer Abbildung, nicht das der Tabellenzelle (siehe \ZSF@imgBudget).
% \makeatletter ist noetig, weil \ZSF@imgBudget ein Kontrollwort mit @ ist;
% \begin{ZSF@box} daneben braucht es nicht — ein Umgebungsname ist Text.
\makeatletter
\NewDocumentEnvironment{figbox}{O{} O{}}
  {\renewcommand{\ZSF@imgBudget}{\ZSFfigMaxHeight}%
   \begin{ZSF@box}{preset/fig}{#1}{#2}}{\end{ZSF@box}} % chktex 9
\makeatother
%
%
%
%
%
%
%
%
%
%
%
%
%
%
\NewDocumentEnvironment{warnbox}{O{Achtung} O{}}
  {\begin{ZSF@box}{zsftitlebox, tone=warn}{#1}{#2}}{\end{ZSF@box}} % chktex 9

%
%
%
%
%
%
%
%
%
%
%
%
%
%
%
%
%
%
%
%
%
%
%
%
%
%
%
%
%
%
%
%
%
%
%
%
\makeatletter
\newcommand{\ZSFgroupcols}{}
\newcommand{\ZSF@groupHead}{}
\newlength{\ZSF@groupOuterAfter}
\pgfkeys{
  /zsf/boxgroup/.is family, /zsf/boxgroup,
  cols/.store in=\ZSFgroupcols,
  head/.store in=\ZSF@groupHead,
  ZSF@groupDefaults/.style={cols=, head=},
}
\NewDocumentEnvironment{ZSFboxgroup}{O{}}{%
  \par
  \setlength{\ZSF@groupOuterAfter}{\ZSFBoxAfterSkip}%
  \addvspace{\ZSFBoxBeforeSkip}%
  \begingroup
  \pgfkeys{/zsf/boxgroup, ZSF@groupDefaults, #1}%
  % Innen keine Abstände: erst dadurch liest sich der Stapel als ein Block.
  \renewcommand{\ZSFBoxBeforeSkip}{0pt}%
  \setlength{\ZSFBoxAfterSkip}{0pt}%
  \ifx\ZSF@groupHead\@empty\else
    \begin{tablebox}%
      \begin{ZSFtable}[zebra=false]{\ZSFgroupcols}%
        \ZSFheaderRow\ZSF@groupHead \\%
      \end{ZSFtable}%
    \end{tablebox}%
  \fi
}{%
  \endgroup
  \par\addvspace{\ZSF@groupOuterAfter}%
}
\makeatother

% ------------------------------------------------------------
% Formula-Box
% ------------------------------------------------------------
% Der Unterschied zur Titelbox ist bewusst und klein: Der Titel sitzt auf dem
% Boxhintergrund statt auf einem eigenen Balken, der Inhalt ist zentriert,
% und zwischen upper und lower liegt eine Trennlinie. Alles andere — Rahmen,
% Ecken, Innenabstand, Titelschrift — kommt aus zsftitlebox.
%
% \ifZSF@inFormulaBox (deklariert in 30_layout_spacing) schaltet
% \textVorBox/\textNachBox auf ihre Formelform um. Das Flag wird in
% 'before upper' gesetzt und ist damit auf die Boxgruppe beschränkt — es
% braucht kein Zurücksetzen.
\makeatletter
\tcbset{preset/formula/.style={
  zsftitlebox,
  align=center,
  ZSF@surface=emphasis,
  ZSF@titlefill=surface,
  frame=hard,
  % Der waagrechte Titelabstand stand hier, weil er in zsftitlebar auf 0pt
  % festgenagelt war und auf der Fläche der Formelbox besonders auffiel. Er
  % kommt jetzt aus `padx` und gilt damit auf jeder Box — der Unterschied der
  % Formelbox ist allein, dass ihr Titel keine eigene Fläche hat.
  bottomtitle=0pt,
  /utils/exec={\renewcommand{\ZSFBoxContext}{\ZSF@inFormulaBoxtrue}},
  after upper={\par},
}}
\NewDocumentEnvironment{formulabox}{O{} O{}}
  {\begin{ZSF@box}{preset/formula}{#1}{#2}}{\end{ZSF@box}} % chktex 9
\makeatother

%
%
%
%
%
%
%
%
%
%
%
%
\newcommand{\ZSFBoxRule}{%
  \noindent\textcolor{FormulaSeparatorColor}{\rule{\linewidth}{\ZSFboxRuleWidth}}%
}

%
%
%
%
%
%
%
%
%
%
%
%
%
%
%
%
%
%
\NewDocumentCommand{\ZSFsep}{o}{%
  \ZSFInterlude{\ZSFspaceSep}%
  \ZSFBoxRule
  \IfValueT{#1}{\ZSFInterlude{\ZSFspaceS}{\ZSFfontBlockLabel #1\par}}%
  \ZSFInterlude{\ZSFspaceSep}%
  \penalty-50%
}

% Dezente Anmerkung als eigene Zeile — für die Erläuterung unter einer Formel
% oder Formelgruppe. Für eine Anmerkung, die auf der Formelzeile bleiben soll,
% \ZSFformulaline; für einen Satz, der eng an die Formel bindet, \textNachBox.
\newcommand{\formulanote}[1]{%
  \ZSFInterlude{\ZSFspaceS}%
  \noindent\ZSFBoxNote{#1}%
  \ZSFInterlude{\ZSFspaceS}%
}
% \ZSFformulaline{<mathe>}{<Anmerkung>} — Formelzeile mit rechtsbündiger
% Kurz-Anmerkung auf derselben Höhe (Gültigkeitsbereich, Einheit, Fallname).
\newcommand{\ZSFformulaline}[2]{%
  \par\noindent$\displaystyle #1$\hfill\ZSFBoxNote{#2}\par
}

% ------------------------------------------------------------
% Runintext + Grafik-Helpers
% ------------------------------------------------------------
% Der Fliesstext-Block nimmt seinen Vorabstand aus DERSELBEN Quelle wie jede
% Box (\ZSF@blockBefore). Vorher stand hier ein fester \ZSFspaceS: Nach einem
% Abschnittsbalken zog eine Box ihren Abstand zusammen, ein runintext nicht —
% derselbe Ort, zwei Abstände, je nachdem was folgt. Sichtbar war das als
% auffällige Lücke zwischen Balken und erstem Satz.
\makeatletter
\newenvironment{runintext}{%
  \ZSF@blockBefore
  \noindent\ZSFReadableAuto\ignorespaces
}{\par}
\makeatother

% Dezente Bildunterschrift (leerer Text -> nichts)
\newcommand{\ZSFfigcaption}[1]{%
  \if\relax\detokenize{#1}\relax\else
    \par\vspace{\ZSFspaceXS}%
    \ZSFBoxNote{#1}\par
  \fi
}

%
%
%
%
%
%
%
%
%
%
%
%
%
%
%
%
%
%
%
%
%
%
\makeatletter
\newcommand{\ZSF@figWidthFrac}{}
\newcommand{\ZSF@figBoxOpts}{}
\newlength{\ZSF@figHeightLimit}
% Oeffnet die figbox mit EXPANDIERTER Optionsliste. Ohne diesen Umweg landet
% der Makroname selbst in der Optionsliste: tcolorbox parst sie, bevor es
% expandiert, und meldet die ganze gesammelte Liste als einen unbekannten
% Schluessel. Der \expandafter am Aufrufort loest genau eine Stufe auf.
\newcommand{\ZSF@figOpen}[2]{\begin{figbox}[{#2}][#1]} % chktex 9
\pgfkeys{
  /zsf/fig/.is family, /zsf/fig,
  width/.store in=\ZSF@figWidthFrac,
  height/.code={\setlength{\ZSF@figHeightLimit}{#1}},
  .unknown/.code={% chktex 26
    \IfBlankTF{#1}%
      {\edef\ZSF@figBoxOpts{\ZSF@figBoxOpts,\pgfkeyscurrentname}}%
      {\edef\ZSF@figBoxOpts{\ZSF@figBoxOpts,\pgfkeyscurrentname=\unexpanded{#1}}}%
  },
  ZSF@figDefaults/.style={width=, height=\ZSFfigMaxHeight},
}

% ------------------------------------------------------------
% \ZSFimage — das Bild OHNE Box
% ------------------------------------------------------------
% \ZSFimage[width=<Anteil>, height=<Mass>]{pfad}
%
% Die Bild-Makros darüber liefern immer einen Block: figbox, Titel, Caption.
% Genau das ist falsch, sobald das Bild INHALT eines anderen Bausteins ist —
% in einer Tabellenzelle (Bauform neben Bezeichnung neben Kennwerten) gibt es
% keinen Platz für eine zweite Box. Bis hierher gab es dafür keinen Weg: rohes
% \includegraphics ist in Kapiteln verboten, \ZSFfig öffnet eine figbox.
%
% Zugleich die EINE Stelle im System, an der \includegraphics steht. Vorher
% dreimal fast gleich (Bildreihe, Einzelbild, Bild-neben-Text) — und wer die
% Skalierung ändern wollte, musste alle drei finden.
%
% Bewusst ohne \par und ohne \centering: Der Baustein, in dem das Bild steht,
% bestimmt Ausrichtung und Umbruch. Beides hier zu setzen machte das Makro in
% einer Tabellenzelle unbrauchbar — dem Fall, für den es existiert.
\newcommand{\ZSF@imgWidthFrac}{1}
\newlength{\ZSF@imgHeightLimit}
% Das Hoehenbudget gehoert dem CONTAINER, nicht dem Bild.
%
% Dasselbe \ZSFimage bedeutet an zwei Orten zwei Groessenordnungen: In einer
% Tabellenzelle muss es in eine Zeile passen (\ZSFimageMaxHeight), als
% Abbildung fuellt es einen Block (\ZSFfigMaxHeight). Das Bild selbst kann das
% nicht wissen — es sieht seinen Container nicht.
%
% Stuende hier ein fester Default, waere genau eine der beiden Lagen still
% falsch: Ein \ZSFimage in einer figbox wuerde ohne Fehler und ohne Warnung auf
% Zellenhoehe gestaucht, und die Hoehe muesste an jedem Aufruf wiederholt
% werden — also genau die Wiederholung, die ein zentraler Default vermeidet.
% Deshalb ist der Default eine INDIREKTION: Jeder Baustein, der Bilder tragen
% kann, bindet sie auf sein eigenes Mass um (figbox/splitbox auf Abbildungs-,
% ZSFtable/valuegrid auf Zellenhoehe). Wer nichts bindet, bekommt die Zelle —
% der Fall, fuer den \ZSFimage existiert.
\newcommand{\ZSF@imgBudget}{\ZSFimageMaxHeight}
\pgfkeys{
  /zsf/image/.is family, /zsf/image,
  width/.store in=\ZSF@imgWidthFrac,
  height/.code={\setlength{\ZSF@imgHeightLimit}{#1}},
  ZSF@imgDefaults/.style={width=1, height=\ZSF@imgBudget},
}
%
% Die senkrechte Luft in einer Tabellenzelle bringt das Bild NICHT selbst mit.
% Sie kommt von \ZSFtableCellPadY (20_tables): »Zellinhalt haelt Abstand zur
% Zeilenkante« gilt fuer ein Bild wie fuer einen Bruch, und der Regler `rows`
% entscheidet darueber. Ein eigener Bildwert waere dieselbe Aussage ein
% zweites Mal.
\newcommand{\ZSFimage}[2][]{%
  \begingroup
  \pgfkeys{/zsf/image, ZSF@imgDefaults, #1}%
  \includegraphics[width=\ZSF@imgWidthFrac\linewidth,
                   height=\ZSF@imgHeightLimit, keepaspectratio]{#2}%
  \endgroup
}

%
%
%
%
%
%
%
\newcounter{ZSF@figNum}
\newlength{\ZSF@figWidth}
\newcommand{\ZSFfig}[4][]{%
  \begingroup
  \renewcommand{\ZSF@figBoxOpts}{}%
  \pgfkeys{/zsf/fig, ZSF@figDefaults, #1}%
  \expandafter\ZSF@figOpen\expandafter{\ZSF@figBoxOpts}{#3}%
    \setcounter{ZSF@figNum}{0}%
    \@for\ZSF@figItem:=#2\do{\stepcounter{ZSF@figNum}}%
    \ifnum\value{ZSF@figNum}>1\relax
      \setlength{\ZSF@figWidth}{\dimexpr
        \ifx\ZSF@figWidthFrac\@empty\ZSFfigUsableFraction\else\ZSF@figWidthFrac\fi
        \linewidth/\value{ZSF@figNum}\relax}%
      \noindent
      \@for\ZSF@figItem:=#2\do{%
        \begin{minipage}[c]{\ZSF@figWidth}%
          \centering
          \ZSFimage[width=1, height=\ZSF@figHeightLimit]{\ZSF@figItem}%
        \end{minipage}%
        \hfill
      }%
      \unskip
      \par
    \else
      \centering
      % Vorbelegung erst hier aufloesen: In der Optionsliste von \ZSFimage
      % stuende sonst eine Fallunterscheidung, die pgfkeys als Wert speichert
      % statt sie auszuwerten.
      \ifx\ZSF@figWidthFrac\@empty\renewcommand{\ZSF@figWidthFrac}{1}\fi
      \ZSFimage[width=\ZSF@figWidthFrac, height=\ZSF@figHeightLimit]{#2}%
    \fi
    \ZSFfigcaption{#4}%
  \end{figbox}% chktex 10
  \endgroup
}

% \ZSFfigside[<Optionen>]{pfad}{Titel}{Inhalt rechts} — Bild links, Inhalt
% rechts. `width` ist hier der Anteil des BILDES (Vorbelegung 0.30); er wird
% erst am Einsatzort gegen die dortige \linewidth ausgewertet.
\newcommand{\ZSFfigside}[4][]{%
  \begingroup
  \renewcommand{\ZSF@figBoxOpts}{}%
  \pgfkeys{/zsf/fig, ZSF@figDefaults, width=0.30, #1}%
  \expandafter\ZSF@figOpen\expandafter{\ZSF@figBoxOpts}{#3}%
    \noindent
    \begin{minipage}[c]{\ZSF@figWidthFrac\linewidth}%
      \centering
      \ZSFimage[width=1, height=\ZSF@figHeightLimit]{#2}%
    \end{minipage}%
    \hfill
    \begin{minipage}[c]{\dimexpr\ZSFfigUsableFraction\linewidth
                              -\ZSF@figWidthFrac\linewidth\relax}%
      \raggedright
      #4%
    \end{minipage}%
    \par
  \end{figbox}% chktex 10
  \endgroup
}
\makeatother

%
%
%
%
%
%
%
%
%
%
%
%
%
\tcbset{split/.style={
  sidebyside,
  sidebyside gap=\ZSFspaceM,
  lower separated=false,
  lefthand width=#1\linewidth,
  before lower={\ZSFBoxSetup},
}}
%
%
%
%
%
%
%
%
%
%
%
%
%
\makeatletter
\newcommand{\ZSF@splitAlign}{top}
\tcbset{
  splitalign/.is choice,
  splitalign/top/.code   ={\renewcommand{\ZSF@splitAlign}{top}},
  splitalign/center/.code={\renewcommand{\ZSF@splitAlign}{center}},
  splitalign/bottom/.code={\renewcommand{\ZSF@splitAlign}{bottom}},
}
\tcbset{ZSF@resolvesplit/.code={\tcbset{sidebyside align=\ZSF@splitAlign}}}
\makeatother
% splitbox ist das rahmenlose, titellose Preset dieses Reglers — der mit
% Abstand häufigste Fall, deshalb ein eigener Name statt vier Optionen.
%
% Sie steht auf derselben Form wie jede andere Box (zsfboxshape) und nimmt mit
% Reglern zurück, was sie rahmenlos macht. Ein 'blankest' wäre kürzer, würde
% aber den Rahmen per Skin entfernen statt per Regler: Ein `frame=soft` vom
% Aufrufer liefe danach ins Leere, weil der Skin gar nichts zeichnet — der
% Regler wäre auf dieser einen Box stillschweigend wirkungslos.
\tcbset{preset/split/.style={
  zsfboxshape,
  bodyparskip=inherit,
  frame=none,
  padx=none, pady=none,
  ZSF@surface=quiet,
  split=0.3,
}}
\makeatletter
\NewDocumentEnvironment{splitbox}{O{}}
  {\renewcommand{\ZSF@imgBudget}{\ZSFfigMaxHeight}%
   \begin{ZSF@box}{preset/split}{}{#1}}{\end{ZSF@box}} % chktex 9
\makeatother

% ------------------------------------------------------------
% Aussagen / Verfahren / Listen
% ------------------------------------------------------------
% Die leise Fassung: dezent, linker Balken im Akzent des Tons. Steht auf derselben
% Form wie die Titelboxen und nimmt genau das zurück, was sie rahmenlos macht
% — Rahmen, Ecken und die obere/untere Polsterung.
%
% Ton, Justierung und Eintritt kommen aus dem Vertrag (zsfboxcontract, via
% zsfboxshape); hier steht nur noch der Unterschied. Die Fläche setzt dieser
% Stil bewusst NICHT — sie kommt vom Flächen-Resolver am Listenende, der Rolle
% und Ton zusammensetzt.
\tcbset{preset/quiet/.style={
  zsfboxshape,
  bodyparskip=inherit,
  frame=none,
  borderline west={\ZSFstatementBarWidth}{0pt}{\ZSFtoneAccent},
  padx=bar, pady=tight,
}}

% ------------------------------------------------------------
% ZSFlist — eine Liste, zwei Nummerierungen
% ------------------------------------------------------------
% Nummerierte Schrittfolge und Faktenliste sind dieselbe Sache: „eine Liste in
% einer optionalen leisen Box". Sie unterscheiden sich in Nummerierung und
% Item-Layout — beides Eigenschaften einer Liste, keine verschiedenen Absichten.
%
%   \begin{ZSFlist}[Titel]            Aufzählung, Marke inline
%   \begin{ZSFlist}[Titel][ordered]   Nummerierung, Marke über dem Text
%   \begin{ZSFlist}                   ohne Titel: ganz ohne Box
%
% Das Item-Layout folgt der Listenart, weil es daran hängt: Ein nummerierter
% Schritt trägt Fliesstext und braucht die Marke als eigene Zeile, ein Fakt
% ist ein Einzeiler. Wer beides entkoppeln muss, bekommt einen Regler — aber
% erst, wenn der Fall wirklich auftritt.
\makeatletter
\newif\ifZSF@listOrdered
\newif\ifZSF@listBoxed
\tcbset{ordered/.is if=ZSF@listOrdered, ordered/.default=true}
\newcommand{\ZSF@listOpen}[1]{%
  \ifZSF@listOrdered
    \begin{enumerate}[
      leftmargin=\ZSFlistLeftMarginOrdered, labelsep=\ZSFlistLabelSep,
      itemsep=\ZSFlistItemSep, topsep=#1, parsep=0pt, partopsep=0pt]%
  \else
    \begin{itemize}[
      leftmargin=\ZSFlistLeftMargin, labelsep=\ZSFlistLabelSep,
      label={\textcolor{\ZSFtoneAccent}{\textbullet}},
      itemsep=\ZSFlistItemSep, topsep=#1, parsep=0pt, partopsep=0pt]%
  \fi
}
\newcommand{\ZSF@listClose}{%
  \ifZSF@listOrdered\end{enumerate}\else\end{itemize}\fi % chktex 9
}
%
%
%
%
%
%
%
%
%
%
\pgfkeys{
  /zsf/listplain/.is family, /zsf/listplain,
  ordered/.code={\tcbset{ordered=#1}},
  ordered/.default=true,
  tone/.code={\tcbset{tone=#1}},
  .unknown/.code={% chktex 26
    \PackageError{ZSF}{%
      Regler '\pgfkeyscurrentname' ist auf einer ZSFlist ohne Titel wirkungslos%
    }{%
      Ohne Titel hat die Liste keine Box, auf die ein Box-Regler wirken kann.
      Entweder einen Titel setzen (dann gelten alle Box-Regler) oder die
      Option weglassen. In dieser Form gelten nur: ordered, tone.%
    }%
  },
}
\NewDocumentEnvironment{ZSFlist}{O{} O{}}
  {\ZSF@listOrderedfalse
   \renewcommand{\ZSF@boxOwnKnobs}{,ordered,}%
   \IfBlankTF{#1}%
     {\ZSF@listBoxedfalse\pgfkeys{/zsf/listplain, #2}\ZSF@listOpen{\ZSFspaceXS}}%
     {\ZSF@listBoxedtrue\begin{defbox}[#1][weight=quiet, #2]\ZSF@listOpen{0pt}}}
  {\ZSF@listClose\ifZSF@listBoxed\end{defbox}\fi} % chktex 9
\makeatother
%
%
%
%
%
%
%
%
%
%
%
%
%
%
\makeatletter
% +m/+g: beide Argumente sind lang. Ein Schritt darf mehrere Absätze tragen
% (Erklärung, Displayformel, Schlusssatz) — xparse-Argumente wären sonst kurz,
% wo das frühere \newcommand lang war.
\NewDocumentCommand{\ZSFItem}{+m +g}{%
  \item
  \IfNoValueTF{#2}%
    {#1}%
    {\IfBlankTF{#1}%
       {#2}%
       {\ifZSF@listOrdered
          {\ZSFfontBlockLabel #1}\\[\ZSFlistMarkGap]#2%
        \else
          {\ZSFfontBlockLabel #1}\,---\,#2%
        \fi}}%
}
\makeatother

%
%
%
%
%
%
%
%
%
%
\makeatletter
\newif\ifZSF@goalStarted
\tcbset{preset/goal/.style={
  zsftitlebox,
  align=center,
  pady=normal,
  frame=strong,
  ZSF@surface=plain,
  /utils/exec={\ZSF@goalStartedfalse},
  after upper={\par},
  atomic,
}}
\makeatother
\NewDocumentEnvironment{goalbox}{O{} O{}}
  {\begin{ZSF@box}{preset/goal}{#1}{#2}}{\end{ZSF@box}} % chktex 9

\newtcbox{\GoalConditionBox}{%
  on line, enhanced,
  boxrule=\ZSFgoalConditionRuleWidth,
  colback=\GoalConditionBackColor,
  colframe=\GoalConditionFrameColor,
  borderline={\ZSFgoalConditionBorderlineWidth}{0pt}{\GoalConditionBorderColor},
  arc=\ZSFgoalConditionInnerArc,
  outer arc=\ZSFboxArc,
  left=\ZSFspaceXS, right=\ZSFspaceXS, top=\ZSFspaceXS, bottom=\ZSFspaceXS,
  nobeforeafter,
  tcbox raise base,
}
\newtcbox{\GoalTargetBox}{%
  on line, enhanced,
  boxrule=\ZSFboxRuleWidth,
  colback=\GoalTargetBackColor,
  colframe=\GoalTargetFrameColor,
  arc=\ZSFboxInnerArc,
  outer arc=\ZSFboxArc,
  left=\ZSFspaceXS, right=\ZSFspaceXS, top=\ZSFspaceXS, bottom=\ZSFspaceXS,
  nobeforeafter,
  tcbox raise base,
}

%
%
%
%
%
%
%
%
%
%
%
%
%
%
\makeatletter
\newcommand{\ZSF@goalLink}[1]{%
  \ifZSF@goalStarted\ZSFInterlude{\ZSFgoalLinkSep}\else\par\ZSF@goalStartedtrue\fi
  #1\par
}
\newcommand{\GoalCondition}[1]{\ZSF@goalLink{\GoalConditionBox{\ZSFfontNote #1}}}
\newcommand{\GoalStep}[1]{\ZSF@goalLink{$\displaystyle #1$}}
\newcommand{\GoalTarget}[1]{\ZSF@goalLink{\GoalTargetBox{$\displaystyle #1$}}}
%
%
%
%
%
%
%
%
%
%
%
\NewDocumentCommand{\ZSFDerivationCase}{s m m m}{%
  \GoalCondition{#2}%
  \IfBooleanTF{#1}%
    {\GoalStep{#3 = \GoalTargetBox{$\displaystyle #4$}}}%
    {\GoalStep{#3}\GoalStep{{}= \GoalTargetBox{$\displaystyle #4$}}}%
}
\makeatother
% Inline-Danger-Tag (rotes Pill-Label für Stolperfallen im Fliesstext)
\newtcbox{\ZSFdanger}{
  colback=DangerTagColor, colframe=DangerTagColor, coltext=\DangerTagTextColor,
  on line, boxsep=\ZSFspaceXS,
  left=\ZSFspaceS, right=\ZSFspaceS, top=\ZSFspaceXS, bottom=\ZSFspaceXS,
  boxrule=0pt, arc=\ZSFboxArc, outer arc=\ZSFboxArc,
  fontupper=\ZSFfontDangerTag
}

% ------------------------------------------------------------
% Valuegrid: kompakte (tabularray-)Wertegrids in einer Panel-Box
% ------------------------------------------------------------
% Interner Panel-Rahmen des Valuegrids — kein öffentlicher Baustein.
% Kapitel greifen ausschliesslich über \begin{valuegrid}{n}[Titel] zu.
\makeatletter
% Die Gitterlinien sind ein Regler, keine Eigenschaft des Bausteins: Ob ein
% Wertegitter seine Zellen einzeln umrandet, hängt am Inhalt und wird pro
% Stelle entschieden (Reglertest, styles/README.md). Ein dichtes Zahlengitter
% will beide Linienrichtungen, eine dreispaltige Zuordnungsliste kommt mit
% waagrechten aus, und wo die Spalten selbst sprechen, stört jedes Lineal.
%
% Das Komma steht IM Wert, nicht in der Spec: `grid=none` liefert sonst einen
% leeren Eintrag in der tblr-Optionsliste.
% Die Linienstaerke ist \ZSFboxRuleWidth und nicht tabularrays Vorgabe: Eine
% Linie ist in diesem System EIN Mass — der Boxrahmen zog 0.5pt, das Gitter
% daneben 0.4pt, und entschieden hatte das niemand.
\newcommand{\ZSF@gridSpec}{{wd=\ZSFboxRuleWidth}}
\newcommand{\ZSF@gridLines}{hlines=\ZSF@gridSpec,vlines=\ZSF@gridSpec,}
\tcbset{
  grid/.is choice,
  grid/both/.code      ={\renewcommand{\ZSF@gridLines}{hlines=\ZSF@gridSpec,vlines=\ZSF@gridSpec,}},
  grid/horizontal/.code={\renewcommand{\ZSF@gridLines}{hlines=\ZSF@gridSpec,}},
  grid/none/.code      ={\renewcommand{\ZSF@gridLines}{}},
  % Dieselben zwei Fragen wie an der ZSFtable, unter denselben Namen und mit
  % derselben Antwort (\ZSF@setRows / \ZSF@setColsep in 20_tables). Sie stehen
  % hier als tcolorbox-Schlüssel, weil das valuegrid eine Box ist und seine
  % Optionsliste die der Box IST — die ZSFtable hat dafür ihre eigene Familie.
  % Zwei Aufrufwege, eine Bedeutung; die Werte kommen aus einer Hand.
  rows/.is choice,
  rows/normal/.code={\ZSF@setRows{normal}},
  rows/roomy/.code ={\ZSF@setRows{roomy}},
  rows/tight/.code ={\ZSF@setRows{tight}},
  colsep/.is choice,
  colsep/normal/.code={\ZSF@setColsep{normal}},
  colsep/tight/.code ={\ZSF@setColsep{tight}},
}
% `rows` und `colsep` gehören in die Vorbelegung: Die Register sind sonst das,
% was die zuletzt gesetzte ZSFtable hinterlassen hat — ausserhalb ihrer Gruppe
% also null, und das Gitter stünde ohne jede Polsterung da.
\tcbset{preset/valuegrid/.style={
  zsftitlebox,
  padx=none, pady=tight,
  ZSF@surface=plain,
  frame=strong,
  grid=both,
  rows=normal, colsep=normal,
  atomic,
}}
% `width` ist Pflicht, nicht Kosmetik: Ohne sie setzt tabularray die Spalten
% auf ihre natuerliche Breite, und das Gitter endet mitten in der Box —
% sichtbar, seit die Optionen ueberhaupt ankommen (siehe unten). Eine Tabelle
% fuellt in diesem System immer die Zeilenbreite (20_tables, Kopf).
%
% Zellpolsterung und Zeilenluft kommen aus DENSELBEN Registern wie bei der
% ZSFtable (\ZSF@setColsep / \ZSF@setRows in 20_tables). Vorher standen hier
% Eigenwerte: `rows=roomy` und `colsep=tight` bedeuteten auf den beiden
% Tabellen-Bausteinen verschiedene Dinge, und \ZSFvaluegridRowSep war ein
% zweites Register fuer dieselbe Frage.
\newcommand{\ZSFvaluegridCommon}{
  width=\linewidth,
  colsep=\ZSF@tblColSep,
  rowsep=\ZSFtableCellPadY,
  \ZSF@gridLines
  rows={valign=m}
}
%
%
%
%
%
%
%
%
%
%
%
%
%
%
%
%
%
%
%
\newcommand{\ZSF@vgOpts}{}
\newcommand{\ZSF@vgOpen}[1]{\begin{tblr}{#1}} % chktex 9
\NewDocumentEnvironment{valuegrid}{O{} O{} m}
  {\renewcommand{\ZSF@boxOwnKnobs}{,grid,rows,colsep,}%
   \renewcommand{\ZSF@imgBudget}{\ZSFimageMaxHeight}%
   \begin{ZSF@box}{preset/valuegrid}{#1}{#2}%
   \edef\ZSF@vgOpts{colspec={X[l] *{#3}{X[c]}},\ZSFvaluegridCommon}%
   \expandafter\ZSF@vgOpen\expandafter{\ZSF@vgOpts}}
  {\end{tblr}\end{ZSF@box}} % chktex 9
\makeatother

%
%
%
%
%
%
%
%
%
%
%
%
%
%
%
%
\newcommand{\ZSFBoxesBreakableOn}{%
  \tcbset{zsftitlebox/.append style={breakable},
          preset/quiet/.append style={breakable}}%
}
\newcommand{\ZSFBoxesBreakableOff}{%
  \tcbset{zsftitlebox/.append style={breakable=false},
          preset/quiet/.append style={breakable=false}}%
}
% Hebt die Platzreserven vor Boxen und Überschriften auf, damit ein
% Anhang dicht packt statt früh umzubrechen.
\newcommand{\ZSFBreakReservesOff}{%
  \renewcommand{\BreakIfNotEnoughSpace}{}%
  \renewcommand{\BreakIfNotEnoughSpaceChapter}{}%
  \renewcommand{\BreakIfNotEnoughSpaceSubsubsection}{}%
}
